\section{Background}

\subsection{Truncated SVD}

For $W \in \mathbb{R}^{m \times n}$:

\begin{equation}
W = U \Sigma V^\top
\end{equation}

The rank-$r$ truncation:

\begin{equation}
W_r = U_{:, :r} \Sigma_{:r, :r} V_{:r, :}^\top
\end{equation}

has Frobenius error $\|W - W_r\|_F = \sqrt{\sum_{i > r} \sigma_i^2}$.

\subsection{Symmetric $\sqrt{\Sigma}$ Split}

Two common factorizations:

\textbf{Asymmetric:}
\begin{equation}
U_r = U_{:, :r}, \quad V_r = \Sigma_{:r, :r} V_{:r, :}^\top
\end{equation}

\textbf{Symmetric:}
\begin{equation}
U_r = U_{:, :r} \Sigma_{:r, :r}^{1/2}, \quad V_r = \Sigma_{:r, :r}^{1/2} V_{:r, :}^\top
\end{equation}

We use the symmetric split for numerical balance: both factors have similar 
magnitudes, reducing FP16 underflow risk.

\subsection{Energy Thresholding}

Given singular values $\sigma_1 \geq \sigma_2 \geq \dots$, the cumulative 
energy fraction at rank $r$ is:

\begin{equation}
E(r) = \frac{\sum_{i=1}^{r} \sigma_i^2}{\sum_{j=1}^{\min(m, n)} \sigma_j^2}
\end{equation}

The \textbf{energy threshold} $\tau \in (0, 1]$ defines the target:

\begin{equation}
r^* = \min\{r : E(r) \geq \tau\}
\end{equation}

This directly bounds the relative Frobenius error:

\begin{equation}
\frac{\|W - W_{r^*}\|_F}{\|W\|_F} = \sqrt{1 - E(r^*)} \leq \sqrt{1 - \tau}
\end{equation}

For $\tau = 0.92$: relative error $\leq \sqrt{0.08} \approx 0.283$.

\subsection{Multi-Head Attention}

For input $X \in \mathbb{R}^{B \times L \times d}$:

\begin{equation}
Q = X W_Q, \quad K = X W_K, \quad V = X W_V
\end{equation}

\begin{equation}
\text{head}_h = \text{softmax}\!\left(\frac{Q_h K_h^\top}{\sqrt{d_h}}\right) V_h
\end{equation}

\begin{equation}
\text{MHSA}(X) = \text{Concat}(\text{head}_1, \dots, \text{head}_H) W_O
\end{equation}

\subsection{T5 Relative Attention Bias}

T5 (Raffel et al., 2020) adds a learned bias to attention scores:

\begin{equation}
\text{scores}_{ij} \leftarrow \text{scores}_{ij} + b_{\text{bucket}(i, j)}
\end{equation}

where $\text{bucket}(i, j)$ maps the relative position to a discrete bin.

Our factorization preserves the relative bias embedding because we only 
replace Q/K/V/O projections, not the bias.

\subsection{Cross-Attention}

Encoder-decoder attention uses \textbf{different} K, V inputs:

\begin{equation}
K = H_\text{enc} W_K, \quad V = H_\text{enc} W_V
\end{equation}

where $H_\text{enc}$ is the encoder's output. Our factorized K, V projections 
apply to any input, so cross-attention works unchanged.

\subsection{Residual Adapters (LoRA)}

We add a trainable low-rank residual:

\begin{equation}
\Delta W = B A
\end{equation}

with $A \in \mathbb{R}^{r_a \times d}$, $B \in \mathbb{R}^{d \times r_a}$, and 
$B = 0$ at initialization.

\subsection{Notation}

\begin{table}[H]
\centering
\small
\begin{tabular}{ll}
\toprule
\textbf{Symbol} & \textbf{Meaning} \\
\midrule
$W$ & Weight matrix \\
$U, \Sigma, V$ & SVD components \\
$\sigma_i$ & $i$-th singular value \\
$r$ & Selected rank \\
$\tau$ & Energy threshold \\
$r_a$ & Residual rank \\
$A, B$ & Adapter matrices \\
$H$ & Number of heads \\
$d_h$ & Head dim \\
\bottomrule
\end{tabular}
\end{table}

